% ============================================================
%  Příprava na maturitu – Mechatronika
%  Kompiluj: pdflatex maturita_mechatronika.tex
%  Kódování: UTF-8
% ============================================================

\documentclass[a4paper,10pt]{article}

% --- Balíčky ---
\usepackage[T1]{fontenc}
\usepackage[utf8]{inputenc}
\usepackage[czech]{babel}
\usepackage[left=2.4cm, right=1.8cm, top=1.8cm, bottom=1.8cm]{geometry}
\usepackage{lmodern}
\usepackage{microtype}
\usepackage{parskip}
\usepackage{titlesec}
\usepackage{enumitem}
\usepackage{xcolor}
\usepackage{hyperref}
\usepackage{tabularx}
\usepackage{booktabs}
\usepackage{array}
\usepackage{tikz}
\usepackage{eso-pic}
\usepackage{mdframed}
\usepackage{amsmath}
\usepackage{amssymb}
\usepackage{pgfplots}
\usetikzlibrary{circuits.logic.US}
\pgfplotsset{compat=1.18}
\usetikzlibrary{arrows.meta, positioning, shapes.geometric, calc}

% --- Barvy ---
\definecolor{accent}{HTML}{03254E}
\definecolor{stripeblue}{HTML}{568EA3}
\definecolor{medgray}{HTML}{888888}
\definecolor{lightblue}{HTML}{EAF2F8}
\definecolor{lightyellow}{HTML}{FDFBE4}
\definecolor{lightgreen}{HTML}{EAF7EA}

% --- Modrý pruh vlevo ---
\AddToShipoutPictureBG{%
  \begin{tikzpicture}[remember picture, overlay]
    \fill[stripeblue]
      (current page.north west)
      rectangle
      ([xshift=10mm] current page.south west);
    \fill[accent!60]
      ([xshift=10mm] current page.north west)
      rectangle
      ([xshift=12mm] current page.south west);
  \end{tikzpicture}%
}

\hypersetup{colorlinks=true, urlcolor=accent, linkcolor=accent}

\titleformat{\section}
  {\large\bfseries\color{accent}}
  {\thesection.}
  {0.5em}
  {}
  [\color{accent}\titlerule]

\titlespacing{\section}{0pt}{16pt}{6pt}

\newmdenv[
  backgroundcolor=lightblue,
  linecolor=stripeblue,
  linewidth=1pt,
  roundcorner=4pt,
  innerleftmargin=8pt,
  innerrightmargin=8pt,
  innertopmargin=6pt,
  innerbottommargin=6pt,
  skipabove=4pt,
  skipbelow=4pt
]{pojmybox}

\newmdenv[
  backgroundcolor=lightyellow,
  linecolor=accent!40,
  linewidth=1pt,
  roundcorner=4pt,
  innerleftmargin=8pt,
  innerrightmargin=8pt,
  innertopmargin=6pt,
  innerbottommargin=6pt,
  skipabove=4pt,
  skipbelow=8pt
]{vysvetlenibox}

\newmdenv[
  backgroundcolor=lightgreen,
  linecolor=accent!30,
  linewidth=1pt,
  roundcorner=4pt,
  innerleftmargin=8pt,
  innerrightmargin=8pt,
  innertopmargin=6pt,
  innerbottommargin=6pt,
  skipabove=4pt,
  skipbelow=8pt
]{vzorcebox}

\pagestyle{plain}

% ============================================================
\begin{document}

% --- Záhlaví ---
\begin{minipage}[t]{0.75\textwidth}
    {\Huge\bfseries\color{accent} Příprava na maturitu}\\[4pt]
    {\large\color{stripeblue} Mechatronika}\\[2pt]
    \textcolor{medgray}{\small Suchánek Lukáš \quad SPŠ Prosek \quad 2026}
\end{minipage}
\hfill
\begin{minipage}[t]{0.22\textwidth}
    \raggedleft
    \begin{tikzpicture}
        \draw[color=stripeblue, rounded corners=4pt, line width=1pt]
            (0,0) rectangle (3,1.2)
            node[midway, align=center, color=accent, font=\small\bfseries]
            {Otázek: 22};
    \end{tikzpicture}
\end{minipage}

\vspace{6pt}
\noindent\color{accent}\rule{\linewidth}{1.5pt}
\vspace{2pt}

% ============================================================
\section{Sekvenční logické řízení, klopné obvody}

\begin{pojmybox}
\textbf{Klíčové pojmy:}
sekvenční obvod, klopný obvod, SR, JK, D, T,
automat (Mealy, Moore), stavový diagram, čítač, registr,
synchronní/asynchronní řízení, hazard
\end{pojmybox}

\begin{vysvetlenibox}
\textbf{Sekvenční obvod}\\
Na rozdíl od kombinačního má sekvenční obvod \textbf{paměť} -- výstup závisí
na vstupech \textit{a} na předchozím stavu. Základním stavebním prvkem je
\textbf{klopný obvod} (bistabilní člen -- může být ve stavu 0 nebo 1).

\textbf{Rozdělení sekvenčních obvodů:}
\begin{itemize}[noitemsep, leftmargin=1.4em, topsep=4pt]
    \item \textbf{Asynchronní} -- stav se mění ihned při změně vstupů;
          rychlejší, ale náchylné k hazardům
    \item \textbf{Synchronní} -- stav se mění pouze při hodinovém impulsu (CLK);
          stabilnější, snadněji se navrhují
\end{itemize}

\textbf{Princip klopného obvodu}\\
Klopný obvod má dva stabilní stavy (0 a 1). Přechod mezi stavy se řídí:
\begin{itemize}[noitemsep, leftmargin=1.4em, topsep=4pt]
    \item \textbf{Úrovní} -- obvod reaguje, když je CLK na vysoké/nízké úrovni
    \item \textbf{Hranou} -- obvod reaguje pouze při náběžné/sestupné hraně CLK
\end{itemize}
Hranové řízení je výhodnější -- zabraňuje vícečetným přepnutím.

\textbf{Přehled klopných obvodů}

\begin{center}
\begin{tabularx}{0.95\linewidth}{@{} l l X l @{}}
    \toprule
    \textbf{Typ} & \textbf{Vstupy} & \textbf{Funkce} & \textbf{Zvláštnost} \\
    \midrule
    RS  & S, R       & S=1: nastav Q=1; R=1: nastav Q=0
                     & Zakázaný stav S=R=1 \\
    D   & D, CLK     & Q přejde na hodnotu D při hraně CLK
                     & Jednoduchý, základ registrů \\
    JK  & J, K, CLK  & J=K=0: paměť; J=1,K=0: set; J=0,K=1: reset;
                       J=K=1: přepnutí
                     & Nemá zakázaný stav \\
    T   & T, CLK     & T=1: přepne stav; T=0: paměť
                     & Základ čítačů \\
    \bottomrule
\end{tabularx}
\end{center}

\textbf{Detailní popis klopných obvodů:}
\begin{itemize}[noitemsep, leftmargin=1.4em, topsep=4pt]
    \item \textbf{RS klopný obvod} -- (set-reset) nejzákladnější typ; S (Set) nastaví Q=1,
          R (Reset) nastaví Q=0. Stav S=R=1 $\rightarrow$ zakázaný stav.
    \item \textbf{D klopný obvod} -- eliminuje zakázaný stav; data (D) se
          zapíší při hraně hodinového signálu (CLK). Základ registrů a pamětí.
    \item \textbf{JK klopný obvod} -- vylepšená verze RS obvodu.
    \begin{center}
    \begin{tabular}{|c|c|l|}
    	\hline
    	J & K & Výtup\\
    	\hline
    	0 & 0 & drží stav\\
    	0 & 1 & vynuluje\\
    	1 & 0 & nastaví\\
    	1 & 1 & překlopí\\
    	\hline
    \end{tabular}
	\end{center}
    \item \textbf{T klopný obvod} -- zjednodušený JK s jedním vstupem;
          při T=1 a CLK=1, se stav překlopí.
\end{itemize}

\begin{vzorcebox}
JK klopný obvod -- charakteristická rovnice:
$Q_{n+1} = J \cdot \overline{Q_n} + \overline{K} \cdot Q_n$\\
D klopný obvod: $Q_{n+1} = D$\\
T klopný obvod: $Q_{n+1} = T \oplus Q_n$\\
SR klopný obvod: $Q_{n+1} = S + \overline{R} \cdot Q_n$ (pro S·R=0)
\end{vzorcebox}

\textbf{Stavový diagram -- vizualizace chování}\\
Stavový diagram znázorňuje přechody mezi stavy pomocí kruhů a šipek:
\begin{itemize}[noitemsep, leftmargin=1.4em, topsep=4pt]
    \item \textbf{Kruh} -- reprezentuje stav (0 nebo 1)
    \item \textbf{Šipka} -- přechod mezi stavy při splnění podmínky
    \item \textbf{Podmínka} -- zápis ve tvaru "vstup/výstup"
\end{itemize}

\textbf{Moorův a Mealyho automat}\\
Automat je model sekvenčního obvodu s definovanými stavy a přechody:

\begin{center}
\begin{tabularx}{0.95\linewidth}{@{} l X X @{}}
    \toprule
    & \textbf{Moorův automat} & \textbf{Mealyho automat} \\
    \midrule
    Výstup závisí na & pouze na stavu & stavu i vstupu \\
    Stabilita výstupu & stabilní, mění se při přechodu & může se měnit i bez přechodu \\
    Počet stavů & obvykle více & obvykle méně \\
    Použití & bezpečnostní systémy & rychlé řídicí automaty \\
    \bottomrule
\end{tabularx}
\end{center}

\textbf{Příklad: Automat pro rozpoznání sekvence 101}\\
\begin{itemize}[noitemsep, leftmargin=1.4em, topsep=4pt]
    \item \textbf{Stavy:} S0 (čekám na 1), S1 (čekám na 0), S2 (čekám na 1), S3 (detekováno)
    \item \textbf{Přechody:} S0--1→S1, S1--0→S2, S2--1→S3, S3--0→S2, S3--1→S1
    \item \textbf{Výstup:} Y=1 pouze ve stavu S3 (Mealyho varianta)
\end{itemize}

\textbf{Čítače a registry}

\begin{itemize}[noitemsep, leftmargin=1.4em, topsep=4pt]
    \item \textbf{Asynchronní čítač} -- výstup předchozího KO je CLK
          následujícího; jednodušší, ale pomalejší (ripple carry).
          \textit{Nevýhoda:} zpoždění se sčítá při průchodu řetězcem.
    \item \textbf{Synchronní čítač} -- všechny KO mají společný CLK;
          rychlejší, bez hazardů. Převodní logika určuje, které KO se přepne.
    \item \textbf{Čítač se zátěží} -- lze přednastavit počáteční hodnotu;
          používá se pro děličky s programovatelným poměrem.
    \item \textbf{Posuvný registr} -- D klopné obvody zapojené do série;
          sériový vstup / výstup dat. Typy:
          \begin{itemize}[noitemsep, leftmargin=1em]
              \item \textbf{SIPO} (Serial-In Parallel-Out) -- sériový vstup, paralelní výstup
              \item \textbf{PISO} (Parallel-In Serial-Out) -- paralelní vstup, sériový výstup
              \item \textbf{SISO} (Serial-In Serial-Out) -- zpoždění signálu
              \item \textbf{PIPO} (Parallel-In Parallel-Out) -- přesun dat mezi registry
          \end{itemize}
\end{itemize}

\textbf{Hazardy v sekvenčních obvodech}\\
Hazard je krátký nežádoucí impuls při změně vstupů:
\begin{itemize}[noitemsep, leftmargin=1.4em, topsep=4pt]
    \item \textbf{Statický hazard} -- výstup má krátce opačnou hodnotu
    \item \textbf{Dynamický hazard} -- výstup několikanásobně přepne
    \item \textbf{Řešení:} synchronní design, redundance v logice, filtry
\end{itemize}
\end{vysvetlenibox}
% ============================================================

% ============================================================
\end{document}
